\documentclass[11pt,a4paper]{article}
\usepackage[utf8]{inputenc}
\usepackage[T1]{fontenc}
\usepackage[spanish]{babel}
\usepackage{hyperref}
\usepackage{geometry}
\geometry{margin=2.5cm}
\usepackage{parskip}

\title{347 el ente en blanco no es de esta era, pero hoy lo parece}
\author{Alejandro Valenzuela Rivera\\\small Universidad Autónoma de Nuevo León}
\date{\today}
\begin{document}
\maketitle

\begin{quote}
\textbf{Sesgo declarado.} Este trabajo se sostiene desde el sesgo declarado de su autor, y no finge neutralidad. Lo que afirma se puede tumbar, y ese es el punto: el falsificador está escrito, no escondido.
\end{quote}

\section*{Resumen}
El ente en blanco no es de una era: hoy parece del estado del arte y mañana no, porque lo que lo sostiene es el acto y no la moda.

\section{Por qué se sostiene}
Un ente que se nombra por su era muere con ella; uno que se nombra por lo que hace sobrevive al que lo puso.

\section{De dónde viene}
La autoridad: 'el ente en blanco no es de esta era, pero hoy lo parece'.

\section{El caso}
El ente en blanco: hoy se parece a la industria, y lo que lo sostiene es más viejo que ella.

\section{Método}
\textquestiondown{}Sigue diciendo lo mismo cuando la industria cambia?

\section{Límites, y qué lo tumba}
Un ente que solo se entiende dentro de su era.

\section{Lo que queda abierto}
\textquestiondown{}El ente se sostiene por su acto, o por parecerse al estado del arte?

\section*{Cómo se leyó esta tesis}
Esto se escribió en el repositorio man-made-empathy, del que este candado es una pieza: candados/01 el ente y el metodo/347 el ente en blanco no es de esta era, pero hoy lo parece.md. La máquina que lo compuso lo hizo desde este mismo archivo, y no agregó nada que no estuviera escrito. Si algo de esto es falso, se dice en el mismo sitio, y se corrige con fecha.

\end{document}
